\documentclass{article}
\usepackage{graphicx}
\usepackage{amsmath}
\usepackage{amsfonts}
\usepackage{bookmark}
\usepackage{booktabs}
\usepackage{exercises}
\graphicspath{ {} }
\usepackage{indentfirst}
\usepackage{amsthm}
\usepackage{yhmath}
\usepackage{geometry}
\usepackage{placeins}
\usepackage{csquotes}
\usepackage{pgfplots}
\newcommand*{\oldneg}{\mathord{\sim}}
\newcommand\showdiv[1]{\overline{\smash{)}#1}}
\newcommand\Mydiv[1]{%
\smash{\raise.3ex\hbox{$\big)$}}\mkern-8mu
        \overline{\enspace\strut#1}}
\setlength{\tabcolsep}{5pt}
\theoremstyle{definition}
\newtheorem{definition}{Definition}

\setlength{\parindent}{.5in}
\setlength{\parskip}{12pt}
\renewcommand{\arraystretch}{2}
\geometry{margin=1in}
\newcommand{\?}{\stackrel{?}{=}}
\pgfplotsset{width=7cm,compat=1.18}\usepgfplotslibrary{statistics}

\title{Chapter Test (pg. 483)}
\begin{document}
\maketitle

\section*{1}
The sequence is arithmetic, with a common difference of 4 and $a_1=5$.

\section*{2}
The sequence is geometric, with a common ratio of 2 and $a_1=3$.

\section*{3}
The sequence is geometric, with a common ratio of $\frac{1}{4}$ and $a_1=40$.

\section*{4}
The sequence is not arithmetic, as $7-4=3$ but $12-7=5$. It is also not geometric,  as $\frac{7}{4}=1.75$ but
$\frac{12}{7}\ne 1.75$.

\section*{5}
$6-1^2=5,6-2^2=2,6-3^2=-3,6-4^2=-10,6-5^2=-19,6-6^2=-30$

\section*{6}
$7\times 1^3=7,7\times 2^3=56,7\times 3^3=189,7\times 4^3=448,7\times 5^3=875,7\times 6^3=1512$

\section*{7}
$4,5\times 4=20,5\times 20=100,5\times 100=500,5\times 500=2500,5\times 2500=12500,5\times 12500=62500$ 

\section*{8}
$-1,-1+6=5,5+6=11,11+6=17,17+6=23,23+6=29$

\section*{9}
The next term is $29$, as the rule is $a_n=5+(n-1)6$.

\section*{10}
The next term is $1875$, as the rule is $a_n=3(5)^{n-1}$.

\section*{11}
The next term is $\frac{10}{25}=\frac{2}{5}$, as the rule is $a_n=\dfrac{6+(n-1)}{5+5(n-1)}$.

\section*{12}
The next term is $8.0$, as the rule is $a_n=1.6+(n-1)1.6$.

\section*{13}
\[ \sum_{i=1}^{48}i=\frac{48(48+1)}{2}=1176 \]

\section*{14}
\[ \sum_{n=1}^{28}n^2=\frac{28(28+1)(2(28)+1)}{6}=7714 \]

\section*{15}
\[ \sum_{i=1}^{10}(4i-9)=\sum_{i=1}^{10}4i+\sum_{i=1}^{10}-9 \]
\[ =4\sum_{i=1}^{10}i-9\sum_{i=1}^{10}1 \]
\[ =4\frac{10(10+1)}{2}-90 \]
\[ =220-90=130 \]

\section*{16}
\[ \sum_{i=1}^{19}(2i+5)=\sum_{i=1}^{19}2i+\sum_{i=1}^{19}5 \]
\[ =2\sum_{i=1}^{19}i+5\sum_{i=1}^{19}1 \]
\[ =2\frac{19(19+1)}{2}+95 \]
\[ =475 \]

\section*{17}
This is a geometric series with $a_1=9$ and $r=2$. Thus, the sum of the first 5 terms is:
\[ 9\frac{1-2^5}{1-2}=279 \]

\section*{18}
This is a geometric series with $a_1=12$ and $r=\frac{1}{3}$. Thus, the sum of the first 6 terms is:
\[ 12\frac{1-(\frac{1}{3})^6}{1-\frac{1}{3}}=\frac{1456}{81} \]

\section*{19}
This is an infinite geometric series with $a_1=8$ and $r=\frac{3}{4}$. Thus, because $|r|<1$, the sum of the series is:
\[ \frac{8}{1-\frac{3}{4}}=32 \]

\section*{20}
This is an infinite geometric series with $a_1=20$ and $r=\frac{3}{10}$. Thus, because $|r|<1$, the sum of the series is:
\[ \frac{20}{1-\frac{3}{10}}=\frac{200}{7} \]

\section*{21}
This repeating decimal can be expressed as the sum of an infinite geometric series with $a_1=0.1$ and $r=0.1$. This has a sum of:
\[ \frac{0.1}{1-0.1}=\frac{1}{9} \] meaning this repeating decimal is $\frac{1}{9}$.

\section*{22}
This repeating decimal can be expressed as the sum of an infinite geometric series with $a_1=0.46$ and $r=0.01$. This has a sum of:
\[ \frac{0.46}{1-0.01}=\frac{0.46}{0.99}=\frac{46}{99} \] meaning this repeating decimal is $\frac{46}{99}$.

\section*{23}
This repeating decimal can be expressed as the sum of an infinite geometric series with $a_1=0.187$ and $r=0.001$. This has a sum of:
\[ \frac{0.187}{1-0.001}=\frac{187}{999} \] meaning this repeating decimal is $\frac{187}{999}$.

\section*{24}
(I was unsure if this answer had $25$ as the repeating part, or $3252525$ as the repeating part, so I gave both answers)\newline
Solution 1: This repeating decimal can be expressed as the sum of $0.3$ and the sum of an infinite geometric series with $a_1=0.025$ and $r=0.01$. This has a sum of:
\[ \frac{0.025}{1-0.01}=\frac{2500}{99000}=\frac{5}{198} \] meaning the fraction is 
\[ 0.3+\frac{5}{198}=\frac{594}{1980}+\frac{50}{1980}=\frac{644}{1980}=\frac{161}{495} \]


Solution 2: This repeating decimal can be expressed as the sum of an infinite geometric series with $a_1=0.3252525$ and $r=0.0000001$. This has a sum of:
\[ \frac{0.3252525}{1-0.0000001}=\frac{3252525}{9999999}=\frac{1084175}{3333333} \]

\section*{25}
The rule is $a_1=2$ and $a_n=a_{n-1}\times 6$.

\section*{26}
The rule is $a_1=3$ and $a_n=a_{n-1}+7$.

\section*{27}
The rule is $a_1=135$ and $a_n=a_{n-1}\times\frac{1}{3}$.

\section*{28}
The rule is $a_1=1$ and $a_n=a_{n-1}\times -3$.

\section*{29}
The first three iterates are $f(x_1)=f(x_0),f(x_2)=f(x_1),f(x_3)=f(x_2)$. Thus
\[ f(x_1)=f(4)=3(4)-7=5,f(x_2)=f(5)=3(5)-7=8,f(x_3)=f(8)=3(8)-7=17 \]

\section*{30}
The first three iterates are $f(x_1)=f(x_0),f(x_2)=f(x_1),f(x_3)=f(x_2)$. Thus
\[ f(x_1)=f(1)=8-5(1)=3,f(x_2)=f(3)=8-5(3)=-7,f(x_3)=f(-7)=8-5(-7)=43 \]

\section*{31}
The first three iterates are $f(x_1)=f(x_0),f(x_2)=f(x_1),f(x_3)=f(x_2)$. Thus
\[ f(x_1)=f(-1)=(-1)^2+2=3,f(x_2)=f(3)=3^2+2=11,f(x_3)=f(11)=11^2+2=123 \]

\section*{32}
(a): $n$ represents the stage of the sequence the quilt is in, and $a_n$ represents how many blue squares are remaining in the quilt. That is, $a_n$ represents how many blue squares are in 
a square checkerboard quilt with a side length of $n$ squares. More specifically, if $n$ is odd, $a_n$ is the number of blue squares in a checkerboard with side length of $n$ formed by 
marking each square on a diagonal starting from an odd row in the first column a blue square. If $n$ is even, it is the number of blue squares when marking each square on from a diagonal 
starting from a sqaure on an even row on the first column.


(b):
\begin{center}
\begin{tabular}{l|l|l|l|l|l|l|l}
$n$ & 1 & 2 & 3 & 4 & 5 & 6 & 7 \\
$a_n$ & 1 & 2 & 5 & 8 & 13 & 18 & 25
\end{tabular}
\end{center}

(c): We can make a table of the results obtained by the given rule:
\begin{center}
\begin{tabular}{l|l|l|l|l|l|l|l}
$n$ & 1 & 2 & 3 & 4 & 5 & 6 & 7 \\
$a_n$ & 1 & 2 & 5 & 8 & 13 & 18 & 25
\end{tabular}
\end{center}
From the table, we can conclude that the results obtained from this rule are identical to the ones in part (b). Thus this sequence is an accurate description of the pattern 
described in part (a).

\section*{33}
We are given that $a_1=3072$ and $r=\frac{1}{4}$. Thus $a_n=3072(\frac{1}{4})^{n-1}$. The values of $n$ that make sense are those for which $a_n$ is a whole number (as we must have a whole number 
numer of actors). Now, because we are dividing
$a_{n-1}$ by $4$ to obtain $a_n$, $a_n$ will stay a whole number as long as $a_{n-1}$ is divisible by four. Hence, we need to determine how many factors of 4 3072 has. 
$3072=4^5\times 3$, meaning it makes sense for there to be 6 rounds, as $a_7=3072(\frac{1}{4})^6=\frac{3}{4}$, which is not a whole number. 

\end{document}
